% GENERATED FILE. Edit canonical lessons, then run npm run generate:books.
% canonical-source-hash: fnv1a64:776b1e30844d8c30
% canonical-lessons: JA-C240-jouzu, JA-C240-heta, JA-C240-suki, JA-C240-kirai, JA-C240-akiru

\chapter{Likes and Skills}
\label{ch:ja-likes-and-skills}
\hlchaptermodality{240}

\begin{chapteropening}
\textbf{By the end of this chapter:} \emph{I can say what I liked and what I was good or bad at.}

Complete the last lesson of chapter 240: I can say what I liked and what I was good or bad at.
\end{chapteropening}

\section[\texorpdfstring{\ja{じょうず} (jōzu)}{じょうず (jōzu)}]{\ja{じょうず} (jōzu) — skilful; good (at)}

\label{lesson:JA-C240-jouzu}



\begin{quote}
Before the new one: say the Japanese for a magic trick, then the Japanese for shogi.
\end{quote}

\subsection*{You'll want to know: \ja{じょうず}}
\textbf{\ja{じょうず}} — \emph{jōzu} — \textquotedblleft{}skilful; good (at)\textquotedblright{}.

\begin{quote}
\textbf{\ja{あねは} \ja{うたが} \ja{じょうずでした。}} — \emph{ane wa uta ga jōzu deshita} — \textquotedblleft{}My elder sister was good at singing.\textquotedblright{}
\end{quote}

\begin{cousinweb}
\textbf{\ja{じょうず}} (\emph{jōzu}) is Sino-Japanese (kango): two Chinese characters read with their on'yomi, readings borrowed from Chinese. They are \emph{jō}, \textquotedblleft{}above\textquotedblright{} (Middle Chinese \emph{dzyangH}) and \emph{zu}, \textquotedblleft{}hand\textquotedblright{} (\emph{syuwX}) (KANJIDIC2; Middle Chinese in Baxter's transcription).
\end{cousinweb}

\begin{grammarlens}[title={What you've built}]
One more word for everyday life at A2.
\end{grammarlens}

\subsection*{Your turn}
\begin{itemize}
\raggedright
  \item \emph{Say it:} \emph{jōzu}
  \item \emph{Say it:} \emph{jōzu}, once more
  \item \emph{Say it:} \emph{ane wa uta ga jōzu deshita} — My elder sister was good at singing.
  \item \emph{Recall:} say the Japanese for a magic trick, then the Japanese for shogi, then say \emph{jōzu} again
\end{itemize}

\begin{tcolorbox}[breakable,colback=teal!4,colframe=teal!35!black,title={Before you move on}]
What does \ja{じょうず} mean? (\textquotedblleft{}Skilful; good (at)\textquotedblright{}.) Say it once more.
\end{tcolorbox}
\section[\texorpdfstring{\ja{へた} (heta)}{へた (heta)}]{\ja{へた} (heta) — unskilful; poor (at)}

\label{lesson:JA-C240-heta}



\begin{quote}
Before the new one: say the Japanese for shogi, then the Japanese for skilful.
\end{quote}

\subsection*{You'll want to know: \ja{へた}}
\textbf{\ja{へた}} — \emph{heta} — \textquotedblleft{}unskilful; poor (at)\textquotedblright{}.

\begin{quote}
\textbf{\ja{おとうとは} \ja{やきゅうが} \ja{へたでした。}} — \emph{otōto wa yakyū ga heta deshita} — \textquotedblleft{}My younger brother was bad at baseball.\textquotedblright{}
\end{quote}

\begin{cousinweb}
\textbf{\ja{へた}} (\emph{heta}) is written with two characters meaning \textquotedblleft{}below\textquotedblright{} and \textquotedblleft{}hand\textquotedblright{}, read as a whole word rather than character by character; its older origin is not settled (JMdict, KANJIDIC2).
\end{cousinweb}

\begin{grammarlens}[title={What you've built}]
One more word for everyday life at A2.
\end{grammarlens}

\subsection*{Your turn}
\begin{itemize}
\raggedright
  \item \emph{Say it:} \emph{heta}
  \item \emph{Say it:} \emph{heta}, once more
  \item \emph{Say it:} \emph{otōto wa yakyū ga heta deshita} — My younger brother was bad at baseball.
  \item \emph{Recall:} say the Japanese for shogi, then the Japanese for skilful, then say \emph{heta} again
\end{itemize}

\begin{tcolorbox}[breakable,colback=teal!4,colframe=teal!35!black,title={Before you move on}]
What does \ja{へた} mean? (\textquotedblleft{}Unskilful; poor (at)\textquotedblright{}.) Say it once more.
\end{tcolorbox}
\section[\texorpdfstring{\ja{すき} (suki)}{すき (suki)}]{\ja{すき} (suki) — liked; fond of}

\label{lesson:JA-C240-suki}



\begin{quote}
Before the new one: say the Japanese for skilful, then the Japanese for unskilful.
\end{quote}

\subsection*{You'll want to know: \ja{すき}}
\textbf{\ja{すき}} — \emph{suki} — \textquotedblleft{}liked; fond of\textquotedblright{}.

\begin{quote}
\textbf{\ja{きょねん} \ja{あねは} \ja{まんがが} \ja{すきでした。}} — \emph{kyonen ane wa manga ga suki deshita} — \textquotedblleft{}Last year my elder sister liked comics.\textquotedblright{}
\end{quote}

\begin{cousinweb}
\textbf{\ja{すき}} (\emph{suki}) is a native Japanese word (wago), inherited rather than borrowed: it comes from \emph{suku}, \textquotedblleft{}to like\textquotedblright{}, and its kanji was attached later as a kun'yomi (JMdict, KANJIDIC2).
\end{cousinweb}

\begin{grammarlens}[title={What you've built}]
One more word for everyday life at A2.
\end{grammarlens}

\subsection*{Your turn}
\begin{itemize}
\raggedright
  \item \emph{Say it:} \emph{suki}
  \item \emph{Say it:} \emph{suki}, once more
  \item \emph{Say it:} \emph{kyonen ane wa manga ga suki deshita} — Last year my elder sister liked comics.
  \item \emph{Recall:} say the Japanese for skilful, then the Japanese for unskilful, then say \emph{suki} again
\end{itemize}

\begin{tcolorbox}[breakable,colback=teal!4,colframe=teal!35!black,title={Before you move on}]
What does \ja{すき} mean? (\textquotedblleft{}Liked; fond of\textquotedblright{}.) Say it once more.
\end{tcolorbox}
\section[\texorpdfstring{\ja{きらい} (kirai)}{きらい (kirai)}]{\ja{きらい} (kirai) — disliked; hated}

\label{lesson:JA-C240-kirai}



\begin{quote}
Before the new one: say the Japanese for unskilful, then the Japanese for liked.
\end{quote}

\subsection*{You'll want to know: \ja{きらい}}
\textbf{\ja{きらい}} — \emph{kirai} — \textquotedblleft{}disliked; hated\textquotedblright{}.

\begin{quote}
\textbf{\ja{おとうとは} \ja{やさいが} \ja{きらいでした。}} — \emph{otōto wa yasai ga kirai deshita} — \textquotedblleft{}My younger brother disliked vegetables.\textquotedblright{}
\end{quote}

\begin{cousinweb}
\textbf{\ja{きらい}} (\emph{kirai}) is a native Japanese word (wago), inherited rather than borrowed: it comes from \emph{kirau}, \textquotedblleft{}to dislike\textquotedblright{}, and it is usually written in kana (JMdict, KANJIDIC2).
\end{cousinweb}

\begin{grammarlens}[title={What you've built}]
One more word for everyday life at A2.
\end{grammarlens}

\subsection*{Your turn}
\begin{itemize}
\raggedright
  \item \emph{Say it:} \emph{kirai}
  \item \emph{Say it:} \emph{kirai}, once more
  \item \emph{Say it:} \emph{otōto wa yasai ga kirai deshita} — My younger brother disliked vegetables.
  \item \emph{Recall:} say the Japanese for unskilful, then the Japanese for liked, then say \emph{kirai} again
\end{itemize}

\begin{tcolorbox}[breakable,colback=teal!4,colframe=teal!35!black,title={Before you move on}]
What does \ja{きらい} mean? (\textquotedblleft{}Disliked; hated\textquotedblright{}.) Say it once more.
\end{tcolorbox}
\section[\texorpdfstring{\ja{あきる} (akiru)}{あきる (akiru)}]{\ja{あきる} (akiru) — to get tired of; to lose interest}

\label{lesson:JA-C240-akiru}



\begin{quote}
Before the new one: say the Japanese for liked, then the Japanese for disliked.
\end{quote}

\subsection*{You'll want to know: \ja{あきる}}
\textbf{\ja{あきる}} — \emph{akiru} — \textquotedblleft{}to get tired of; to lose interest\textquotedblright{}.

\begin{quote}
\textbf{\ja{いもうとは} \ja{にんぎょうに} \ja{あきました。}} — \emph{imōto wa ningyō ni akimashita} — \textquotedblleft{}My younger sister got tired of her dolls.\textquotedblright{}
\end{quote}

Polite \textbf{\ja{あきます}}, \emph{akimasu}; an ichidan verb.

\begin{cousinweb}
\textbf{\ja{あきる}} (\emph{akiru}) is a native Japanese word (wago), inherited rather than borrowed from Chinese; its kanji was attached later as a kun'yomi (JMdict, KANJIDIC2).
\end{cousinweb}

\begin{grammarlens}[title={What you've built}]
One more word for everyday life at A2.
\end{grammarlens}

\subsection*{Your turn}
\begin{itemize}
\raggedright
  \item \emph{Say it:} \emph{akiru}
  \item \emph{Say it:} \emph{akiru}, once more
  \item \emph{Say it:} \emph{imōto wa ningyō ni akimashita} — My younger sister got tired of her dolls.
  \item \emph{Recall:} say the Japanese for liked, then the Japanese for disliked, then say \emph{akiru} again
\end{itemize}

\begin{tcolorbox}[breakable,colback=teal!4,colframe=teal!35!black,title={Before you move on}]
What does \ja{あきる} mean? (\textquotedblleft{}To get tired of; to lose interest\textquotedblright{}.) Say it once more.
\end{tcolorbox}
